\documentclass[preview]{standalone}

\usepackage{amsmath}
\usepackage{amssymb}
\usepackage{stellar}
\usepackage{definitions}

\begin{document}

\id{psicologia-comunicazione-non-verbale}
\genpage

\section{Comunicazione non verbale}

\begin{snippet}{comunicazione-non-verbale-introduzione}
    Per gran parte dell'evoluzione umana, la comunicazione è avvenuta attraverso
    azioni, movimenti corporei, gesti e sguardo. Su questa piattaforma non
    verbale si è radicata e sviluppata la facoltà del linguaggio.
\end{snippet}

\begin{snippet}{asimmetria-verbale-nonverbale}
    Esiste un'asimmetria strutturale tra linguaggio verbale e sistemi non
    verbali. Il linguaggio verbale richiede sempre un supporto non verbale o
    extralinguistico; i sistemi non verbali, invece, possono funzionare in modo
    autonomo nel proprio ambito di riferimento.
\end{snippet}

\begin{snippetexample}{mimica-facciale-example}{Mimica facciale}
    La \emph{mimica facciale} è un insieme coerente e articolato di espressioni
    visibili, capace di mostrare segnali emotivi, intenzionali e relazionali di
    diversa natura.
\end{snippetexample}

\begin{snippet}{digitale-analogico}
    Il linguaggio verbale ha valore digitale, perché dipende dalla presenza o
    assenza di tratti convenzionali e arbitrari. I sistemi non verbali hanno
    invece valore motivato, iconico e analogico, perché mantengono variazioni
    continue e graduate collegate alla realtà che evocano.
\end{snippet}

\begin{snippetdefinition}{silenzio-definition}{Silenzio}
    Il \emph{silenzio} è una forma di comunicazione non verbale fondata
    sull'assenza di parola. Il suo valore comunicativo dipende dalla sua
    ambiguità e dalle regole culturali che stabiliscono dove, quando, come e
    perché usarlo.
\end{snippetdefinition}

\begin{snippet}{silenzio-culture}
    Nelle culture occidentali il silenzio è spesso interpretato come minaccia o
    mancanza di cooperazione nella conversazione. In alcune culture orientali,
    invece, può segnalare fiducia, confidenza e armonia.
\end{snippet}

\begin{snippetdefinition}{teoria-programmi-affettivi-comunicazione-definition}{Teoria dei programmi affettivi}
    La \emph{teoria dei programmi affettivi}, sostenuta da Ekman, afferma che
    esiste un'universalità delle espressioni facciali specifiche per ciascuna
    emozione di base.
\end{snippetdefinition}

\begin{snippetdefinition}{regole-esibizione-definition}{Regole di esibizione}
    Le \emph{regole di esibizione} sono regole attraverso cui le persone
    modulano volontariamente l'espressione delle proprie emozioni in funzione
    della situazione.
\end{snippetdefinition}

\begin{snippet}{regole-esibizione-tipi}
    Le principali regole di esibizione sono accentuazione, attenuazione,
    neutralizzazione e simulazione. Esse mostrano che l'espressione facciale non
    dipende solo dallo stato emotivo interno, ma anche dal contesto e dalle
    norme sociali.
\end{snippet}

\begin{snippetdefinition}{prospettiva-comunicativa-espressioni-definition}{Prospettiva comunicativa}
    La \emph{prospettiva comunicativa} sostiene che le espressioni facciali
    abbiano valore sociale, poiché permettono di comunicare agli altri, in modo
    flessibile, i propri obiettivi e desideri.
\end{snippetdefinition}

\begin{snippetdefinition}{effetto-audience-definition}{Effetto audience}
    L'\emph{effetto audience} è l'influenza della presenza reale o implicita di
    altri individui sulla comparsa, sul mantenimento e sul cambiamento delle
    espressioni facciali.
\end{snippetdefinition}

\begin{snippetdefinition}{gesto-definition}{Gesto}
    Un \emph{gesto} è un'azione motoria che partecipa all'elaborazione del
    significato di un messaggio.
\end{snippetdefinition}

\begin{snippet}{tipi-gesti}
    I principali tipi di gesti sono gesti iconici, pantomime, emblemi, gesti
    deittici, gesti motori e lingua dei segni. Gesti e parole sono
    interdipendenti: sono generati dalla stessa rappresentazione mentale,
    manifestano la stessa intenzione comunicativa e vengono pianificati in
    sintonia con il contesto.
\end{snippet}

\begin{snippetdefinition}{prossemica-definition}{Prossemica}
    La \emph{prossemica} riguarda la percezione, l'organizzazione e l'uso dello
    spazio, della distanza e del territorio nei confronti degli altri.
\end{snippetdefinition}

\begin{snippetdefinition}{territorialita-definition}{Territorialità}
    La \emph{territorialità} è la regolazione dello spazio personale e sociale
    attraverso cui gli individui definiscono distanza, vicinanza e confini nelle
    relazioni.
\end{snippetdefinition}

\begin{snippet}{zone-prossemiche}
    La gestione dello spazio personale può essere distinta in zona intima, zona
    personale, zona sociale e zona pubblica. Le norme che regolano queste
    distanze variano tra culture della vicinanza e culture della distanza.
\end{snippet}

\begin{snippetdefinition}{aptica-definition}{Aptica}
    L'\emph{aptica} è il sistema di significazione e segnalazione che riguarda
    le azioni di contatto corporeo nei confronti degli altri.
\end{snippetdefinition}

\begin{snippet}{aptica-culture}
    Il contatto corporeo è un atto comunicativo non verbale primario e risponde
    a un bisogno fondamentale della specie umana. Esistono però differenze
    culturali: alcune culture valorizzano il contatto, altre lo limitano o lo
    regolano in modo più restrittivo.
\end{snippet}

\end{document}
